\section*{Erd\H{o}s Problem 650: matching integers to distinct multiples}

\paragraph{Statement and answer.}
Let \(f(m)\) be the largest universally guaranteed number of distinct
elements of an \(m\)-element set \(A\subseteq\{1,\ldots,N\}\) that
can be matched to distinct divisible integers in any interval of
length \(2N\). Then
\[
 \boxed{f(m)=\min\{m,\lceil2\sqrt m\rceil\}.}
\]
We prove the lower bound even for open intervals, and give a matching
upper construction whose endpoints are not multiples of any member
of \(A\). Thus endpoint conventions do not affect the answer.
In particular the displayed question \(f(m)\leq\sqrt m\) on the
problem page is not the correct exact bound.

\paragraph{Attribution.}
This reconstructs the corrected proof of van Doorn, Li, and Tang,
\emph{Optimal bounds for an Erd\H{o}s problem on matching integers
to distinct multiples}, \url{https://arxiv.org/abs/2603.28636},
Theorems 3.1 and 4.1. That paper attributes the original strategy to
GPT-5.4 Pro and the repair of the exceptional injection to Aristotle.
The present argument is an attributed reconstruction, not a new
discovery or a claim of a new Lean verification.
See \url{https://www.erdosproblems.com/650}.

\paragraph{The matching reduction.}
Let the right vertices be the integers in an open interval
\(I=(x,x+2N)\), and join \(a\in A\) to its multiples in \(I\).
For \(S\subseteq A\), write \(\Gamma(S)\) for its neighbourhood.
If every nonempty \(S\) satisfies
\[
 |\Gamma(S)|\geq2\sqrt{|S|},                                      \tag{1}
\]
Hall's theorem supplies a matching of size at least
\(\min(m,\lceil2\sqrt m\rceil)\).
Indeed, set
\(\delta=\max_{S\subseteq A}(|S|-|\Gamma(S)|)\geq0\).
Adding \(\delta\) universal right vertices gives Hall's condition,
and a matching of all \(m\) left vertices loses at most \(\delta\)
edges when the added vertices are deleted.
Since \(s-2\sqrt s\) increases for integer \(s\geq1\), (1) yields
\(\delta\leq\max(0,m-2\sqrt m)\), proving the assertion.

It suffices to treat \(N=\max A\): any longer interval contains
an open subinterval of the required smaller length.
If \(x\notin N\mathbb Z\), split the interval at \(b_0=x+N\), with
the midpoint in the left half. For each \(a\in S\), let \(u_a\)
be the largest multiple of \(a\) at most \(b_0\). Then
\[
 x<u_a\leq b_0<u_a+a<x+2N.
\]
The final inequality is strict: equality would require \(a=N\)
and \(x\in N\mathbb Z\).
The map \(a\mapsto(u_a,u_a+a)\) into the two halves of \(\Gamma(S)\)
is injective because the difference recovers \(a\).
Their sizes therefore have product at least \(|S|\), and the
arithmetic--geometric mean inequality proves (1).

\paragraph{The exceptional endpoint case.}
Suppose \(x\in N\mathbb Z\), so \(b_0=x+N\) is an integer.
Match \(N\) to \(b_0\), and delete these two vertices.
Let \(A'=A\setminus\{N\}\). For a subset \(S\subseteq A'\) we
construct an injection into its remaining left-half and right-half
neighbours. Define \(u_a\) as before, and set
\[
 \phi(a)=
 \begin{cases}
 (u_a,u_a+a),&
   a\nmid b_0 \quad\text{(type 1)},\\
 (b_0-2a,b_0+2a),&
   a\mid b_0,\ 2a<N,\ 2a\nmid b_0 \quad\text{(type 2)},\\
 (b_0-a,b_0+a),&
   a\mid b_0,\ [\,2a\geq N\ \text{or}\ 2a\mid b_0\,]
          \quad\text{(type 3)}.
 \end{cases}                                                     \tag{2}
\]
Every coordinate is a multiple of \(a\), belongs to the correct
open half, and differs from \(b_0\).
Within each type, the coordinate difference recovers \(a\).
For mixed types, suppose two images coincide.
Their midpoint is then \(b_0\).
For types 1 and 2, with arguments \(a,a'\), the difference gives
\(a=4a'\), whence \(b_0=u_a+2a'\) is divisible by \(2a'\),
a contradiction.
For types 1 and 3, the difference gives \(a=2a'<N\).
Type 3 therefore requires \(2a'\mid b_0\), whereas
\(b_0=u_a+a'\) is an odd multiple of \(a'\), again impossible.
For types 2 and 3 the difference gives \(a'=2a\);
type 3 implies \(a'\mid b_0\), contradicting the type 2 condition.
Thus (2) is injective.

It follows that (1) holds in the reduced graph.
For \(m\geq2\), adding back the reserved edge gives at least
\[
 1+\min(m-1,2\sqrt{m-1})\geq\min(m,2\sqrt m)
\]
matched vertices; integrality gives the desired ceiling.
The case \(m=1\) is immediate. This completes the lower bound,
including the endpoint case where the simpler injection fails.

\paragraph{A rectangular upper construction.}
We show that a set of \(st\) integers can have all its multiples
in an interval confined to just \(s+t\) integers.
Cases with \(\min(s,t)=1\) satisfy the upper bound trivially,
so assume \(2\leq s\leq t\).
Let \(D=\operatorname{lcm}(1,\ldots,st)\), and let \(\mathcal P\)
consist of the primes \(p>st\) dividing at least one nonzero integer
\[
 q+rD,\qquad 0<|q|<s,\quad |r|<t.
\]
This is a finite set. For each \(p\in\mathcal P\), avoid the at most
\(st<p\) residue classes \(-i-jD\), \(1\leq i\leq s\),
\(1\leq j\leq t\). The ordinary Chinese remainder theorem lets us
choose \(M>2s+2tD\) avoiding all of them simultaneously. Set
\[
 a_{ij}=M+i+jD,\qquad
 A=\{a_{ij}:1\leq i\leq s,\ 1\leq j\leq t\}.
\]
These \(st\) integers are distinct because \(D>s\).

We claim
\[
 \gcd(a_{ij},a_{k\ell})\mid i-k.                                \tag{3}
\]
There is nothing to prove if \(i=k\). Otherwise put
\(q=i-k\ne0\), \(r=j-\ell\).
Any common prime divisor larger than \(st\) belongs to
\(\mathcal P\), contrary to the choice of \(M\).
For a common prime \(p\leq st\), we have \(v_p(D)>v_p(q)\):
if \(p\nmid q\), this follows from \(p\mid D\); otherwise
\(q^2<s^2\leq st\) implies \(q^2\mid D\).
Thus \(v_p(q+rD)=v_p(q)\).
Since the gcd divides \(q+rD\), (3) follows prime by prime.

The generalized Chinese remainder theorem now supplies \(x_0\)
with \(x_0\equiv i\pmod{a_{ij}}\) for every \(i,j\).
Adding a large common multiple makes \(x=x_0-M>1\).
Put \(N=\max A=M+s+tD\) and \(I=(x,x+2N)\).
For each \(a_{ij}\), \(x_0-i\) is a multiple, and the preceding
and second following multiples satisfy
\[
 (x_0-i)-a_{ij}=x-2i-jD<x,
\]
\[
 (x_0-i)+2a_{ij}=x+3M+i+2jD>x+2N.
\]
Consequently the only possible multiples in \(I\) are
\[
 x_0-i\quad\hbox{and}\quad x_0+M+jD.
\]
Both lie strictly inside \(I\); the strict preceding and following
inequalities also exclude multiples at its endpoints.
There are at most \(s+t\) integers in this union, so no divisibility
matching can be larger.

\paragraph{All values of \(m\).}
Set \(q=\lceil2\sqrt m\rceil\),
\(s=\lfloor q/2\rfloor\), \(t=\lceil q/2\rceil\).
Then \(st=\lfloor q^2/4\rfloor\geq m\) and \(s+t=q\).
If \(q\geq m\), the trivial upper bound \(f(m)\leq m\) suffices.
Otherwise \(s,t\geq2\), and the rectangular construction applies.
Keep any \(m\) members of its set, including its largest member.
The interval retains length \(2\max A\), while its union of
available multiples cannot grow. Hence \(f(m)\leq q\).
Together with the lower bound, this proves the exact formula.

\paragraph{Completion estimate.}
\textbf{COMPLETION: 100\%}
